% Standalone candidate proof dossier for the near-worst bootstrap/interface cluster.
\documentclass[../main.tex]{subfiles}
\begin{document}

% === SOLUTION HEADER =========================================================
% Certification is not recorded here: research/program/ledger.yaml owns the
% proofs[] mode and review path, and the review's front matter owns identities.
%
% ledger-node : lem:half; lem:whitening; thm:bootstrap; lem:crude; cor:loglog; prop:ceiling
% refines : lem:half; lem:whitening; thm:bootstrap; lem:crude; cor:loglog;
% author : /root/kls_bootstrap_author
% date : 2026-08-25
% ============================================================================
\section*{Solution: the near-worst bootstrap and its covariance interface}
\label{sec:sol-kls-bootstrap-interface}

\paragraph{Setup.}
For a log-concave probability measure $\nu$, let
\[
  I_\nu(p)=\inf\{\nu^+(E):\nu(E)=p\},\qquad
  h_\nu=\inf_E\frac{\nu^+(E)}{\min(\nu(E),1-\nu(E))},
\]
and let $\hstar_n$ be the infimum of $h_\nu$ over isotropic log-concave laws on
$\mathbb R^n$.  Along stochastic localization write
\[
  p_t=\mu_t(E),\quad q_t=1-p_t,\quad s_t=p_tq_t,
  \quad A_t=\operatorname{Cov}(\mu_t),
\]
\[
  \bar e_t(E)=\mu_t^+(E)-h_{\mu_t}\min(p_t,q_t),qquad
  X_t=(\lambda_{\max}(A_t)-1)_+,qquad
  \Xi_T=\int_0^T\mathbb E X_t\,\dd t.
\]
For $0<\eta<1/2$ set
$\tau_\eta=\inf\{t:|p_t-1/2|>\eta\}$.  The mass martingale satisfies
\begin{equation}
  \dd[p]_t=s_t r_t\,\dd t,
  \qquad s_t r_t\leq\frac14\lambda_{\max}(A_t)
  \leq\frac14(1+X_t).
  \label{eq:sol-mass-qv-bound}
\end{equation}
All stopping and expectation arguments below are first made in the compact smooth class and
then passed through the manuscript's uniform approximation convention.  Since $0\leq p_t\leq1$,
the stopped mass martingales are uniformly integrable; no unbounded optional-stopping theorem is
being invoked.

\subsection*{1. Two deterministic isoperimetric comparisons}

\begin{lemma}[Balance pins the Cheeger ratio; Lemma~\ref{lem:half}]
\label{lem:sol-half}
For every log-concave $\nu$,
\[
  h_\nu=2I_\nu(1/2).
\]
If $\nu(E)=1/2$, then the profile excess and Cheeger-line excess agree:
\[
  \nu^+(E)-I_\nu(1/2)
  =\nu^+(E)-h_\nu\min(\nu(E),1-\nu(E)).
\]
\end{lemma}

\begin{proof}
The generalized isoperimetric profile of a log-concave measure is symmetric,
$I_\nu(p)=I_\nu(1-p)$, is concave on $(0,1)$ in the standard generalized sense, and has
$I_\nu(0+)=0$ \cite{Bobkov1999LogConcave,Milman2009Isoperimetric}.  Concavity at the origin
implies that $p\mapsto I_\nu(p)/p$ is nonincreasing on $(0,1/2]$: if $0<p<q\leq1/2$, then
$I_\nu(p)\geq(p/q)I_\nu(q)$.  Therefore
\[
  \inf_{0<p\leq1/2}\frac{I_\nu(p)}p=2I_\nu(1/2).
\]
Symmetry gives the same infimum on $[1/2,1)$ after division by $1-p$, proving the first
identity.  The second is immediate because $\min(1/2,1/2)=1/2$.
\end{proof}

\begin{lemma}[Whitening comparison; Lemma~\ref{lem:whitening}]
\label{lem:sol-whitening}
Let $\nu$ be log-concave on $\mathbb R^n$ with nondegenerate covariance $A_\nu$.  Then
\begin{equation}
  h_\nu\geq\frac{\hstar_n}{\lambda_{\max}(A_\nu)^{1/2}}.
  \label{eq:sol-whitening}
\end{equation}
In particular, almost surely
$h_{\mu_t}\geq\hstar_n\lambda_{\max}(A_t)^{-1/2}$ for $t>0$.  Moreover,
$\hstar_n$ is nonincreasing in $n$.
\end{lemma}

\begin{proof}
Let $a_\nu$ be the centroid, $T=A_\nu^{1/2}$, and
$\widetilde\nu=(A_\nu^{-1/2})_\#(\nu-a_\nu)$.  Then $\widetilde\nu$ is isotropic and
log-concave, so $h_{\widetilde\nu}\geq\hstar_n$, and
$\nu=T_\#\widetilde\nu+a_\nu$.  Put $L=\|T\|_{\mathrm{op}}$.
For every Borel set $S$ and every $\delta>0$,
\[
  T^{-1}(S_\delta-a_\nu)
  \supseteq\bigl(T^{-1}(S-a_\nu)\bigr)_{\delta/L}.
\]
After subtracting the common mass, dividing by $\delta$, and taking the lower limit,
\[
  \nu^+(S)\geq L^{-1}\widetilde\nu^+(T^{-1}(S-a_\nu))
  \geq\frac{\hstar_n}{L}\min(\nu(S),1-\nu(S)).
\]
Taking the infimum in $S$ and using $L=\lambda_{\max}(A_\nu)^{1/2}$ proves
\eqref{eq:sol-whitening}.  An isotropic starting law is full-dimensional, and the localization
density is positive on the same support, so $A_t$ remains nondegenerate at finite $t$.

For monotonicity, if $m\leq n$ and $\rho$ is isotropic log-concave on $\mathbb R^m$, then
$\rho\otimes\gamma_{n-m}$ is isotropic log-concave on $\mathbb R^n$.  Cylinder competitors
give $h_{\rho\otimes\gamma_{n-m}}\leq h_\rho$.  Taking first the infimum over all $\rho$ and
then using the definition in dimension $n$ yields $\hstar_n\leq\hstar_m$.
\end{proof}

\subsection*{2. The bootstrap comparison}

\begin{theorem}[Bootstrap comparison; Theorem~\ref{thm:bootstrap}]
\label{thm:sol-bootstrap}
Let $n\geq2$, $\varepsilon\in(0,1]$, and let $\mu$ be isotropic log-concave on
$\mathbb R^n$ with
$h_\mu\leq(1+\varepsilon)\hstar_n$.  Let $E$ be any measurable set with
$\mu(E)=1/2$ and
$e_0=e_0(E)=\bar e_0(E)$, and fix $\eta\in(0,1/2)$.  For every $t>0$,
\begin{equation}
  \mathbb E[\bar e_t(E)\mathbf 1_{\{t<\tau_\eta\}}]
  \leq e_0+h_\mu\left(
    \frac\varepsilon2+\eta+\frac12\mathbb P(\tau_\eta\leq t)
    +\frac14\mathbb E X_t\right),
  \label{eq:sol-bootstrap-main}
\end{equation}
\begin{equation}
  \mathbb P(\tau_\eta\leq t)
  \leq\frac1{4\eta^2}\left(t+\int_0^t\mathbb E X_s\,\dd s\right).
  \label{eq:sol-bootstrap-exit}
\end{equation}
Consequently,
\begin{equation}
\begin{split}
  \int_0^T\mathbb E[\bar e_t(E)\mathbf 1_{\{t<\tau_\eta\}}]\,\dd t
  \leq{}&Te_0+h_\mu\left[
    \left(\frac\varepsilon2+\eta\right)T+
    \frac{T^2}{16\eta^2}+
    \frac{T\Xi_T}{8\eta^2}+\frac{\Xi_T}{4}\right].
\end{split}
\label{eq:sol-bootstrap-integrated}
\end{equation}
If $0<T<1/8$, $\eta=T^{1/3}$, and $\varepsilon\leq T^{1/3}$, then for a universal $C$,
\begin{equation}
  \int_0^T\mathbb E[\bar e_t(E)\mathbf 1_{\{t<\tau_\eta\}}]\,\dd t
  \leq Te_0+C h_\mu(T^{4/3}+\Xi_T).
  \label{eq:sol-bootstrap-clean}
\end{equation}
\end{theorem}

\begin{proof}
Write $A=\{t<\tau_\eta\}$ and $P=\mathbb P(A^c)$.  The perimeter supermartingale gives
\begin{equation}
  \mathbb E[\mu_t^+(E)\mathbf 1_A]
  \leq\mathbb E\mu_t^+(E)
  \leq\mu^+(E)=\frac{h_\mu}{2}+e_0,
  \label{eq:sol-bootstrap-perimeter}
\end{equation}
where the last equality is Lemma~\ref{lem:sol-half}.  This uses a deterministic time $t$;
there is no optional-stopping step for the perimeter process.

On $A$, $\min(p_t,q_t)\geq1/2-\eta$.  Lemma~\ref{lem:sol-whitening} and the elementary
inequality
\begin{equation}
  \lambda^{-1/2}\geq1-\frac12(\lambda-1)_+,
  \qquad \lambda>0,
  \label{eq:sol-sqrt-tangent}
\end{equation}
give
\begin{equation}
\begin{split}
  \mathbb E[h_{\mu_t}\min(p_t,q_t)\mathbf 1_A]
  &\geq\hstar_n(1/2-\eta)
     \mathbb E[\lambda_{\max}(A_t)^{-1/2}\mathbf 1_A]\\
  &\geq\hstar_n(1/2-\eta)(1-P-\tfrac12\mathbb E X_t).
\end{split}
\label{eq:sol-bootstrap-profile-raw}
\end{equation}
The last lower bound may be negative.  The manuscript's abbreviated proof silently multiplied
it by a lower bound for $\hstar_n/h_\mu$, which is only legitimate when its bracket is
nonnegative.  We now record the required two-case argument.

Set $Y=\mathbb E X_t$, $a=1/2-\eta$, and
$\rho=\hstar_n/h_\mu$.  Near-worstness and the definition of $\hstar_n$ give
\begin{equation}
  \frac1{1+\varepsilon}\leq\rho\leq1,
  \qquad\text{hence}\qquad \rho\geq1-\varepsilon.
  \label{eq:sol-rho-range}
\end{equation}
If $1-P-Y/2\geq0$, combine
\eqref{eq:sol-bootstrap-perimeter}--\eqref{eq:sol-rho-range} to obtain
\[
\begin{split}
 \mathbb E[\bar e_t\mathbf 1_A]-e_0
 &\leq h_\mu\{1/2-(1-\varepsilon)a(1-P-Y/2)\}\\
 &\leq h_\mu(\eta+P/2+Y/4+\varepsilon/2).
\end{split}
\]
Indeed, $1/2-a(1-P)=\eta+aP\leq\eta+P/2$, the $Y$ contribution is at most $Y/4$, and the
$\varepsilon$ contribution is at most $\varepsilon/2$.

If $1-P-Y/2<0$, use only the nonnegativity of
$h_{\mu_t}\min(p_t,q_t)\mathbf 1_A$ in
\eqref{eq:sol-bootstrap-perimeter}.  The case assumption gives
$P/2+Y/4>1/2$, and therefore
\[
  \mathbb E[\bar e_t\mathbf 1_A]-e_0
  \leq h_\mu/2
  <h_\mu(\eta+P/2+Y/4+\varepsilon/2).
\]
Thus \eqref{eq:sol-bootstrap-main} holds in both cases.

For the exit estimate, stop the bounded continuous martingale
$M_u=p_{u\wedge\tau_\eta}-1/2$.  On $\{\tau_\eta\leq t\}$, continuity gives
$\sup_{u\leq t}|M_u|\geq\eta$.  The $L^2$ maximal inequality, the martingale isometry for the
bounded stopping time, and \eqref{eq:sol-mass-qv-bound} yield
\[
\begin{split}
  \mathbb P(\tau_\eta\leq t)
  &\leq\eta^{-2}\mathbb E M_t^2
   =\eta^{-2}\mathbb E[p]_{t\wedge\tau_\eta}\\
  &\leq\frac1{4\eta^2}\int_0^t
       \mathbb E[\mathbf 1_{\{s<\tau_\eta\}}(1+X_s)]\,\dd s\\
  &\leq\frac1{4\eta^2}\left(t+\int_0^t\mathbb E X_s\,\dd s\right).
\end{split}
\]
This proves \eqref{eq:sol-bootstrap-exit}.  All integrands are nonnegative, so Tonelli gives
\[
\begin{split}
  \int_0^T\mathbb P(\tau_\eta\leq t)\,\dd t
  &\leq\frac1{4\eta^2}\left\{
      \frac{T^2}{2}+\int_0^T(T-s)\mathbb E X_s\,\dd s\right\}\\
  &\leq\frac1{4\eta^2}\left(\frac{T^2}{2}+T\Xi_T\right).
\end{split}
\]
Integrating \eqref{eq:sol-bootstrap-main} and inserting this estimate gives exactly
\eqref{eq:sol-bootstrap-integrated}, including the constants $1/16$, $1/8$, and $1/4$.

Finally, $T<1/8$ ensures $\eta=T^{1/3}<1/2$.  With
$\varepsilon\leq T^{1/3}$,
\[
  (\varepsilon/2+\eta)T\leq\tfrac32T^{4/3},\qquad
  \frac{T^2}{16\eta^2}=\frac1{16}T^{4/3},\qquad
  \frac{T\Xi_T}{8\eta^2}=\frac18T^{1/3}\Xi_T\leq\frac18\Xi_T.
\]
Together with the remaining $\Xi_T/4$, this proves
\eqref{eq:sol-bootstrap-clean}, for example with $C=2$.
\end{proof}

\subsection*{3. Evaluating the interface}

\begin{lemma}[Crude evaluation; Lemma~\ref{lem:crude}]
\label{lem:sol-crude}
For every isotropic log-concave $\mu$ on $\mathbb R^n$ and every $1/n\leq T\leq1$,
\[
  \Xi_T(\mu)\leq1+\log(nT).
\]
\end{lemma}

\begin{proof}
Taking the trace in the covariance SDE gives
\[
  \dd\operatorname{Tr}A_t=\dd N_t-\operatorname{Tr}(A_t^2)\,\dd t
\]
for a continuous local martingale $N$.  After localization, expectation kills $N$ and the
drift is nonpositive.  Since $\operatorname{Tr}A_t\geq0$, Fatou on removing the localization
gives
$\mathbb E\operatorname{Tr}A_t\leq\operatorname{Tr}A_0=n$.  Hence
$\mathbb E X_t\leq\mathbb E\lambda_{\max}(A_t)\leq n$.

For $t>0$, the Brascamp--Lieb covariance cap $A_t\preceq t^{-1}I$ gives
$X_t\leq(t^{-1}-1)_+\leq t^{-1}$.  Therefore
\[
  \mathbb E X_t\leq\min(n,t^{-1}),
\]
and, because $T\geq1/n$,
\[
  \Xi_T\leq\int_0^{1/n}n\,\dd t+
             \int_{1/n}^T\frac{\dd t}{t}
  =1+\log(nT).
\]
This is only the crude upper fence.  In accordance with
\texttt{obs:crude-insufficient}, it is not used as a dimension-free KLS input.
\end{proof}

\begin{corollary}[Polylogarithmic evaluation; Corollary~\ref{cor:loglog}]
\label{cor:sol-loglog}
Under Assumption~\ref{hyp:KI}, for
$t_1(n)=c_0(\log n)^{-C_2}$,
$\mathbb E\|A_t\|_{\mathrm{op}}\leq C_1$ on $[0,t_1(n)]$.  If
$t_1(n)\leq T\leq1$, then
\begin{equation}
  \Xi_T(\mu)\leq C_1t_1(n)+\log(T/t_1(n))
  \leq C(1+\log\log n).
  \label{eq:sol-loglog-interface}
\end{equation}
If also $T<1/8$, $\varepsilon\leq T^{1/3}$, and $\mu,E$ satisfy the hypotheses of
Theorem~\ref{thm:sol-bootstrap}, then
\begin{equation}
  \int_0^T\mathbb E[\bar e_t(E)\mathbf 1_{\{t<\tau_\eta\}}]\,\dd t
  \leq Te_0+C h_\mu(T^{4/3}+1+\log\log n),
  \qquad \eta=T^{1/3}.
  \label{eq:sol-loglog-supply}
\end{equation}
\end{corollary}

\begin{proof}
On $[0,t_1]$, $X_t\leq\|A_t\|_{\mathrm{op}}$, so the assumed expectation bound contributes
at most $C_1t_1$.  On $[t_1,T]$, Brascamp--Lieb gives $X_t\leq t^{-1}$.  Integration proves
the first inequality in \eqref{eq:sol-loglog-interface}.  Since $T\leq1$ and
$t_1=c_0(\log n)^{-C_2}$,
\[
  \log(T/t_1)\leq C_2\log\log n-\log c_0,
\]
and the fixed low-dimensional cases can be absorbed into the universal constant.  This proves
the second inequality.  Substitution into \eqref{eq:sol-bootstrap-clean} proves
\eqref{eq:sol-loglog-supply}.  The published discharge
Corollary~\ref{cor:KI-discharged} establishes Assumption~\ref{hyp:KI} with $C_2=2$; no numerical
evidence is used here.
\end{proof}

\subsection*{4. The all-measure relative-scale ceiling}

\begin{proposition}[Relative-scale ceiling; Proposition~\ref{prop:ceiling}]
\label{prop:sol-ceiling}
Fix $\kappa\in(0,1]$ and a universal $T_0>0$ such that
$9(1+\kappa)T_0\leq1/2$.  If
\begin{equation}
  \Xi_{T_0}(\mu)\leq\kappa T_0
  \label{eq:sol-ceiling-assumption}
\end{equation}
held for every isotropic log-concave $\mu$, then KLS would follow directly, with no geometric
input.  Thus using the clean bootstrap upper bound to certify a relative-scale propagation
estimate would demand an all-measure covariance input already sufficient for KLS; no converse
or equivalence is asserted.
\end{proposition}

\begin{proof}
For every positive semidefinite $A$,
$\lambda_{\max}(A)\leq1+(\lambda_{\max}(A)-1)_+$.  Hence
\eqref{eq:sol-ceiling-assumption} gives
\begin{equation}
  \int_0^{T_0}\mathbb E\lambda_{\max}(A_t)\,\dd t
  \leq(1+\kappa)T_0.
  \label{eq:sol-ceiling-lambda}
\end{equation}
Fix a balanced cut, $p_0=1/2$, and let
$\tau=\inf\{t:p_t\notin[1/3,2/3]\}$.  The stopped mass martingale is bounded.  If
$\tau\leq T_0$, its displacement reaches $1/6$, so the $L^2$ maximal inequality and the
martingale isometry give
\[
\begin{split}
  \mathbb P(\tau\leq T_0)
  &\leq36\,\mathbb E[p]_{T_0\wedge\tau}\\
  &\leq9\int_0^{T_0}\mathbb E[\mathbf 1_{\{t<\tau\}}
                       \lambda_{\max}(A_t)]\,\dd t\\
  &\leq9(1+\kappa)T_0\leq\frac12,
\end{split}
\]
where the second line is \eqref{eq:sol-mass-qv-bound} without the last replacement by $X_t$,
and the third is \eqref{eq:sol-ceiling-lambda}.  Therefore, with probability at least $1/2$,
$\min(p_{T_0},q_{T_0})\geq1/3$.  The balanced-survival lemma
(Lemma~\ref{lem:survival-implies-kls}) now supplies a universal lower bound for the perimeter
of every balanced cut and hence proves KLS.

Finally, \eqref{eq:sol-bootstrap-clean} bounds the bootstrap error by
$C h_\mu(T^{4/3}+\Xi_T)$.  Making this particular certificate no larger than
$\kappa h_\mu T$ at a sufficiently small universal time asks, term by term, for
$\Xi_T\lesssim\kappa T$ (as well as $T^{1/3}\lesssim\kappa$).  The implication just proved
shows why such an \emph{all-measure} input is already KLS-strength.  This is the precise
method-specific meaning of the proposition's final sentence; it is not a logical necessity for
every possible propagation proof.
\end{proof}

\paragraph{Obstructions respected.}
Theorem~\ref{thm:sol-bootstrap} respects \texttt{obs:circularity}: its profile lower bound comes
from the external worst-case constant $\hstar_n$ and the explicit near-worst assumption
$h_\mu\leq(1+\varepsilon)\hstar_n$, not from an assumed lower bound on the random localized
profile.  Lemma~\ref{lem:sol-crude} respects \texttt{obs:crude-insufficient} by proving and
labeling the $\log n$ estimate only as an insufficient fence; the actual corollary uses the
published small-time covariance input.  Proposition~\ref{prop:sol-ceiling} proves only the
sufficient implication that generates \texttt{obs:relative-ceiling}; it does not claim an
equivalence and does not confuse the all-measure condition with the near-worst route.

\end{document}
